\section{Scope, contributions, and non-claims}
\label{sec:scope}

This paper develops a \emph{deterministic substrate theory fragment}.
Its purpose is not to present a finished replacement for established physics, but to put the central mathematical ingredients of the proposal into a cleaner and more testable form.
The main contribution is the assembly of one coherent chain---from lattice postulate to continuum equations, from geometric diagnostics to the localized-mode program---so that the computational burden and the remaining conceptual gaps become explicit rather than implicit.

\subsection{Positioning relative to established physics}
The Standard Model and general relativity are treated here as \emph{empirically validated effective frameworks}.
Nothing in this manuscript is intended to diminish their established experimental domain of validity.
The present work asks a different question:
whether some of their structural features can be reinterpreted as emergent regularities of a deeper deterministic medium.

\subsection{Core contributions}
The theory fragment advanced here is organized around five specific claims.
\begin{enumerate}
\item \textbf{A continuum substrate model with explicit ontology.}
The fundamental variable is an embedding field $\vecX: \Omega\times\R \to \R^4$ for an intrinsic three-dimensional brane, with external time as evolution parameter.
The continuum dynamics is written as a local hyperelastic action built from the induced metric and an isotropic reference metric.

\item \textbf{A lattice realization of the same substrate.}
A cubic lattice embedded in $\R^4$ is treated as the candidate microphysical realization.
The continuum model is its long-wavelength effective description.
This makes anisotropy, Brillouin-zone structure, and axis-aligned carrier mixing part of the physical hypothesis rather than mere numerical artifacts.

\item \textbf{A geometric contraction channel.}
Because transverse deformation changes intrinsic distances, prestretch induces a transverse-to-lateral backreaction.
In a weak and slowly varying regime this yields a candidate mechanism for a gravity-like contraction field sourced by stored substrate energy.

\item \textbf{Geometric-phase gauge diagnostics.}
An ordered narrowband carrier sector supports Berry or Wilczek--Zee transport.
The corresponding connection and curvature are interpreted as the natural geometric variables of an emergent gauge description.
On the cubic lattice, the axis-aligned triplet $\Psi=(\psi_x,\psi_y,\psi_z)^\top$ provides the minimal non-Abelian testbed.

\item \textbf{Localized modes as the central validation problem.}
Particle-like behavior is assigned to nonlinear localized modes of the substrate.
The manuscript formulates the corresponding linearized eigenproblems, clarifies the role of nonlinear self-guidance, and identifies proton- and neutron-like baryon-scale triplet states as the most informative first simulation targets.
\end{enumerate}

\subsection{Immediate computational objective}
The paper is written to support a concrete next step:
a numerical search for stable three-channel localized modes on the cubic lattice whose size, energy, branch content, and holonomy diagnostics can be compared against baryon-scale targets.
This is narrower and more credible than claiming a direct derivation of QCD.
It is also more aligned with the actual structure of the current substrate hypothesis, because the axis-aligned triplet is intrinsically a three-channel object.

\subsection{Non-claims}
To keep the scope honest, the present version does \emph{not} claim:
\begin{itemize}
\item a derivation of the full Standard Model gauge group, generations, or measured coupling constants,
\item a quantitative reduction of the contraction channel to Einstein gravity or a derived value of $G$,
\item a derivation of full QCD from first principles,
\item a complete scattering theory or multi-particle theory,
\item an already completed proton, neutron, or electron simulation.
\end{itemize}
Instead, the claim is that the mathematical core of the proposal can be stated more cleanly than before and that the next decisive step is computational rather than rhetorical.
